\section{Volterra integral equations: a primer}
\label{app:volterra}

This appendix is self-contained and assumes no prior exposure to integral
equations. Everything the main text needs about Volterra equations is here.

\subsection{What is an integral equation, and why ``Volterra''?}

An \emph{integral equation} has the unknown function under an integral sign,
just as a \emph{differential} equation has it under a derivative. The kind we meet
is the \emph{linear Volterra equation of the second kind}:
\begin{equation}
\xi(t)=s(t)+\int_0^t K(t,u)\,\xi(u)\dd u,\qquad t\in[0,T].
\label{eq:volterra2}
\end{equation}
Here $s$ (the ``forcing'' or ``free term'') and $K$ (the ``kernel'') are given, and
$\xi$ is unknown. ``Second kind'' means the unknown also appears \emph{outside} the
integral (the lone $\xi(t)$ on the left); a \emph{first-kind} equation
$s(t)=\int_0^t K(t,u)\xi(u)\dd u$ has the unknown only inside and is much more
delicate (it is a deconvolution and is typically ill-posed).

The name \emph{Volterra} refers to the upper limit being the running time $t$:
the integral only looks at the \emph{past} $u\le t$. Contrast a \emph{Fredholm}
equation, whose integral runs over a fixed full interval $\int_0^T$. The Volterra
(causal) structure is what makes \eqref{eq:volterra2} so well-behaved: it is, in a
precise sense, ``lower-triangular in time.''

\begin{intuition}
Think of \eqref{eq:volterra2} as a feedback rule unrolling in time: the value
$\xi(t)$ ``now'' equals an external input $s(t)$ \emph{plus} a weighted echo of
everything $\xi(u)$ that happened before, with weight $K(t,u)$. Because the echo
only reaches back into the past, you can in principle sweep forward and build
$\xi$ left to right --- nothing in the future is needed to determine the present.
That causal, triangular character is the source of all the good news below.
\end{intuition}

\subsection{The operator form and the Neumann (resolvent) series}

Define the \emph{Volterra operator}
\begin{equation}
(\mathcal K\xi)(t):=\int_0^t K(t,u)\,\xi(u)\dd u,
\end{equation}
so \eqref{eq:volterra2} is $\xi=s+\mathcal K\xi$, i.e.\ $(\mathcal I-\mathcal
K)\xi=s$. Formally inverting the operator as a geometric series,
\begin{equation}
\xi=(\mathcal I-\mathcal K)^{-1}s=\sum_{n=0}^{\infty}\mathcal K^n s
= s+\mathcal K s+\mathcal K^2 s+\cdots
\label{eq:neumann}
\end{equation}
This is the \emph{Neumann series}. The operator $\mathcal K^n$ has its own kernel,
the \emph{iterated kernel} $K_n(t,u)$ (with $K_1=K$ and
$K_{n+1}(t,u)=\int_u^t K(t,r)K_n(r,u)\dd r$), and the sum
$R(t,u)=\sum_{n\ge1}K_n(t,u)$ is the \emph{resolvent kernel}. With it,
\begin{equation}
\xi(t)=s(t)+\int_0^t R(t,u)\,s(u)\dd u,
\label{eq:resolvent}
\end{equation}
so solving the equation amounts to computing $R$ \emph{once} and then a single
integral against any forcing $s$. $R$ satisfies the resolvent equations
$R=K+K\!\circ\!R=K+R\!\circ\!K$, where $\circ$ is the Volterra composition above.

\begin{intuition}
$\mathcal K s$ is the first-generation echo of the input, $\mathcal K^2 s$ the
echo of the echo, and so on; \eqref{eq:neumann} simply sums all generations. The
resolvent $R$ packages every generation into one kernel --- it is the
``total-feedback'' Green's function of the problem. (In the Hawkes setting $R$ is
literally the renewal density of the branching tree, \S\ref{sec:no-covariates}.)
\end{intuition}

\subsection{Existence and uniqueness: why Volterra is so forgiving}

\begin{theorem}[Existence and uniqueness on a finite horizon]\label{thm:vol-exist}
Suppose $K$ is bounded, $|K(t,u)|\le \bar K$ on the triangle $0\le u\le t\le T$,
and $s\in C[0,T]$. Then \eqref{eq:volterra2} has a unique continuous solution on
$[0,T]$, namely the (uniformly convergent) Neumann series \eqref{eq:neumann}.
\end{theorem}
\begin{proof}
By induction the iterated operator satisfies
$|(\mathcal K^n\xi)(t)|\le \bar K^{\,n}\,\|\xi\|_\infty\,t^n/n!$: indeed
$|(\mathcal K\xi)(t)|\le \bar K\|\xi\|_\infty t$, and
$|(\mathcal K^{n+1}\xi)(t)|\le \int_0^t \bar K\,|(\mathcal K^n\xi)(r)|\dd r
\le \bar K^{\,n+1}\|\xi\|_\infty\int_0^t r^n/n!\dd r=\bar K^{\,n+1}\|\xi\|_\infty t^{n+1}/(n+1)!$.
Hence $\|\mathcal K^n\|_{C[0,T]}\le(\bar K T)^n/n!\to 0$, the series
$\sum_n\mathcal K^n$ converges in operator norm, $(\mathcal I-\mathcal K)$ is
boundedly invertible, and the fixed point \eqref{eq:neumann} is unique.
\end{proof}

\begin{intuition}
The decisive estimate is the factorial $\bar K^n t^n/n!$. Each time the echo
passes through the kernel it must integrate over a \emph{shrinking} triangle of
``time still available,'' and $n$ such nested integrations contribute a $1/n!$ that
overwhelms any fixed gain $\bar K^n$. So on a finite horizon the loop is
\emph{always} invertible --- there is no smallness condition on the kernel at all.
This is the single most important structural fact, and it is why the MBPP equation
is solvable even when the excitation is covariate-modulated and time-varying
(\S\ref{sec:excitation-cov}): nothing in the proof used convolution structure.
\end{intuition}

\noindent
A smallness condition \emph{does} appear if one wants the resolvent globally on
$[0,\infty)$: for a convolution kernel $K(t,u)=\phi(t-u)$, Young's inequality gives
$\|\phi^{\conv n}\|_1\le\|\phi\|_1^n$, so $R=\sum_n\phi^{\conv n}\in L^1(0,\infty)$
iff $\|\phi\|_1=\kappa<1$ --- exactly subcriticality.

\subsection{Three solution strategies, by kernel type}

The structure of $K$ dictates the best method. The main text uses all three.

\paragraph{(i) Convolution kernel $K(t,u)=\phi(t-u)$ $\Rightarrow$ transforms / LTI.}
Then $\mathcal K\xi=\phi\conv\xi$ and the system is time-invariant. Laplace or
Fourier turn the convolution into a product (Appendix~\ref{app:transforms}):
$\hat\xi=\hat s/(1-\hat\phi)$. The resolvent is itself a convolution,
$R(t,u)=h(t-u)$ with $h=\sum_n\phi^{\conv n}$, and for exponential $\phi$ it is a
single exponential. \emph{This is \S\ref{sec:no-covariates}--\ref{sec:baseline-cov}.}

\paragraph{(ii) Separable (degenerate) kernel
$K(t,u)=\sum_{r=1}^{\rho}a_r(t)\,b_r(u)$ $\Rightarrow$ finite ODE.}
A finite-rank kernel collapses the integral equation onto $\rho$ ordinary
differential equations. Indeed, defining the states $y_r(t)=\int_0^t b_r(u)\xi(u)\dd u$,
the equation \eqref{eq:volterra2} becomes $\xi(t)=s(t)+\sum_r a_r(t)\,y_r(t)$, and
differentiating each state gives the closed ODE system
\begin{equation}
y_r'(t)=b_r(t)\,\xi(t)=b_r(t)\Bigl(s(t)+\textstyle\sum_{r'}a_{r'}(t)y_{r'}(t)\Bigr),
\qquad y_r(0)=0.
\label{eq:sep-ode}
\end{equation}
If the $a_r,b_r$ are constants (or pure exponentials), \eqref{eq:sep-ode} has
\emph{constant} coefficients --- a linear time-invariant ODE; if the $a_r$ carry an
extra $t$-dependence (as the covariate factor $e^{\delta^\top Z(t)}$ does), it is a
\emph{time-varying} linear ODE. \emph{An exponential triggering shape is separable
of rank one per component pair, which is precisely why
\S\ref{sec:excitation-cov} reduces to the LTV system \eqref{eq:ltv}.}

\paragraph{(iii) General kernel $\Rightarrow$ numerical quadrature (Nyström).}
When $K$ is neither convolution nor separable, discretise the integral on a grid
$0=u_0<\cdots<u_N=T$. Replacing $\int_0^{t_i}$ by a quadrature rule turns
\eqref{eq:volterra2} into a \emph{lower-triangular} linear system
\begin{equation}
\xi_i = s_i + \sum_{k\le i} w_{ik}\,K(t_i,u_k)\,\xi_k,\qquad i=0,\dots,N,
\end{equation}
which is solved by forward substitution in $O(N^2)$ (the diagonal term is handled
by isolating $\xi_i$). Convergence is at the order of the quadrature rule. This is
the fallback for, e.g., a power-law shape under covariate-modulated excitation.

\begin{intuition}
The three strategies are the same idea --- invert $\mathcal I-\mathcal K$ ---
specialised to how much structure $K$ has. Maximum structure (convolution): a
formula. Medium structure (finite rank): a handful of ODEs. No structure: march
forward in time filling in $\xi$ one step at a time, which the triangular
(causal) form always permits.
\end{intuition}

\begin{remark}[First-kind equations and why we avoid them]
The compensator's \emph{inverse} problem --- recovering $\phi$ from observed mean
behaviour --- is a first-kind (deconvolution) Volterra problem and is ill-posed
(small data errors blow up). The main text never inverts in this direction for
estimation; it always poses a well-posed second-kind forward solve inside a
likelihood, which is far more stable.
\end{remark}
